\section{Giới thiệu đề tài}

\subsection{Bối cảnh và vấn đề}

Trong bối cảnh thương mại điện tử và bán lẻ đa quốc gia phát triển mạnh mẽ, các doanh nghiệp phân phối hàng hóa toàn cầu phải vận hành đồng thời trên hàng trăm thị trường quốc gia với danh mục sản phẩm lên đến hàng chục nghìn mặt hàng (Stock Keeping Units -- SKU). Hoạt động giao dịch hàng ngày tạo ra một khối lượng lớn dữ liệu vận hành phát sinh từ các hệ thống xử lý giao dịch trực tuyến (Online Transaction Processing -- OLTP), bao gồm thông tin đơn đặt hàng, lịch trình giao vận, phân khúc khách hàng, cấu trúc giá bán, chính sách chiết khấu và lợi nhuận ròng trên từng dòng sản phẩm. Tuy nhiên, việc lưu trữ dữ liệu dưới dạng các tệp tin phẳng rời rạc hoặc trong các cơ sở dữ liệu giao dịch chưa qua chuẩn hóa khiến doanh nghiệp gặp nhiều rào cản khi cần tổng hợp báo cáo phân tích lịch sử dài hạn và hỗ trợ ra quyết định ở cấp độ chiến lược.

Qua khảo sát bài toán thực tế tại doanh nghiệp bán lẻ toàn cầu Global Superstore với hơn 51.290 bản ghi giao dịch chi tiết trải rộng trên 147 quốc gia trong giai đoạn 2011--2014, nhóm thực hiện nhận diện ba vấn đề cốt lõi cần giải quyết. Thứ nhất, dữ liệu nguồn xuất phát từ các hệ thống vận hành có độ phân mảnh và không đồng nhất cao về mặt định dạng, điển hình là sự pha trộn giữa nhiều quy chuẩn ngày tháng, tình trạng khuyết thiếu mã bưu chính (Postal Code) tại các đơn hàng ngoài lãnh thổ Hoa Kỳ và sự trùng lặp thông tin danh mục khách hàng cũng như địa lý qua hàng chục nghìn dòng giao dịch. Nếu truy vấn trực tiếp trên nguồn dữ liệu thô này, chi phí tính toán sẽ rất lớn và dễ dẫn đến sai lệch số liệu tài chính khi tổng hợp theo nhiều chiều.

Thứ hai, trong công tác quản trị hiệu quả kinh doanh, ban điều hành đối mặt với bài toán suy giảm biên lợi nhuận cục bộ do chính sách chiết khấu thiếu kiểm soát. Mặc dù tổng doanh thu toàn cầu tăng trưởng đều đặn qua từng năm và đạt ngưỡng 12,64 triệu USD, nhiều khu vực thị trường (như Trung Đông -- EMEA hoặc một số quốc gia tại Châu Phi và Đông Nam Á) và một số phân nhóm ngành hàng (điển hình là nhóm Bàn làm việc -- Tables thuộc ngành Nội thất) lại ghi nhận mức lợi nhuận ròng âm nặng nề khi áp dụng mức chiết khấu từ 30\% đến trên 50\%. Các báo cáo thống kê tĩnh truyền thống chỉ dừng lại ở việc mô tả con số tổng hợp mà không cung cấp khả năng khoan sâu (Drill-down) hoặc xoay chiều (Pivot) linh hoạt để bóc tách nguyên nhân gây thất thoát dòng tiền.

Thứ ba, các nhà quản lý thiếu một hệ thống hỗ trợ quyết định (Decision Support System -- DSS) mang tính định lượng để dự báo nhu cầu thị trường và đánh giá rủi ro trước khi ban hành chính sách. Khi đứng trước các câu hỏi quản trị như cần chuẩn bị ngân sách và lượng hàng tồn kho bao nhiêu cho hai quý tiếp theo, nên ghép cặp những sản phẩm nào vào cùng một gói khuyến mãi bán chéo (Cross-selling) để gia tăng giá trị giỏ hàng trung bình, hay việc điều chỉnh tăng/giảm 5\% tỷ lệ chiết khấu sẽ tác động như thế nào đến tổng lợi nhuận dựa trên độ co giãn của cầu theo giá, việc ra quyết định dựa trên cảm tính kinh nghiệm thường tiềm ẩn rủi ro tài chính cao. Những vấn đề thực tiễn này đặt ra yêu cầu cấp thiết về việc xây dựng một kho dữ liệu (Data Warehouse) tập trung, chuẩn hóa theo mô hình đa chiều kết hợp cùng các mô hình phân tích dự báo và khai phá dữ liệu tự động.

\subsection{Mục tiêu thực hiện}

Đề tài hướng tới mục tiêu thiết kế, xây dựng và kiểm chứng một hệ thống Kho dữ liệu (Data Warehouse) tích hợp Hệ hỗ trợ quyết định (Decision Support System) hoàn chỉnh từ khâu tiếp nhận dữ liệu thô đến khâu trực quan hóa và khuyến nghị hành động cho bài toán Global Superstore. Ở tầng nền tảng dữ liệu, hệ thống cần thiết lập một quy trình trích xuất, biến đổi và nạp dữ liệu (Extract -- Transform -- Load, viết tắt là ETL) tự động theo kiến trúc phân tầng Medallion (Bronze -- Silver -- Gold) trên hệ quản trị cơ sở dữ liệu quan hệ PostgreSQL. Quy trình này phải đảm bảo tính lũy đẳng (Idempotency), khả năng lưu vết kiểm toán từng đợt nạp thông qua mã băm nội dung SHA-256, khả năng nạp gia tăng (Incremental Loading) đối với các lô dữ liệu giao dịch mới phát sinh mà không cần xử lý lại dữ liệu lịch sử, đồng thời bảo toàn tuyệt đối độ chính xác của các chỉ số tài chính giữa các tầng dữ liệu.

Ở tầng lưu trữ và truy vấn phân tích trực tuyến (Online Analytical Processing -- OLAP), mục tiêu của đề tài là chuyển đổi cấu trúc dữ liệu bảng phẳng ban đầu thành mô hình hình sao (Star Schema) tuân thủ phương pháp luận của Ralph Kimball. Mô hình cần tách biệt rõ ràng bảng sự kiện trung tâm (\texttt{gold.fact\_sales}) lưu trữ các độ đo định lượng với năm bảng chiều (\texttt{dim\_date}, \texttt{dim\_customer}, \texttt{dim\_product}, \texttt{dim\_geography}, \texttt{dim\_ship\_mode}) sử dụng khóa thay thế số nguyên (Surrogate Keys). Trên nền tảng mô hình hình sao này, hệ thống phải hiện thực hóa đầy đủ các phép toán OLAP cốt lõi (Roll-up, Drill-down, Slice, Dice, Pivot, CUBE và các hàm cửa sổ phân tích) cùng hệ thống các bảng tổng hợp vật lý (Materialized Views) nhằm tối ưu hóa thời gian phản hồi truy vấn phục vụ bảng điều khiển quản trị.

Ở tầng hỗ trợ ra quyết định (DSS), mục tiêu trọng tâm là phát triển ba mô-đun phân tích chuyên sâu có khả năng vận hành trực tiếp trên dữ liệu sạch từ tầng Gold và lưu trữ kết quả đầu ra vào lược đồ chuyên biệt (\texttt{dss}). Mô-đun thứ nhất thực hiện phân tích dự báo chuỗi thời gian (Predictive Analytics) để ước lượng doanh số bán hàng và sản lượng tiêu thụ trong tương lai kèm theo khoảng tin cậy thống kê 95\%, đồng thời đánh giá khách quan chất lượng dự báo qua các chỉ số sai số thực nghiệm MAE, RMSE và MAPE. Mô-đun thứ hai ứng dụng khai phá luật kết hợp (Association Rule Mining) trên hàng chục nghìn giỏ hàng thực tế để phát hiện các mối tương quan mua sắm đồng thời giữa các mặt hàng và tiểu ngành hàng dựa trên các độ đo Support, Confidence và Lift. Mô-đun thứ ba xây dựng bộ máy mô phỏng kịch bản What-If dựa trên mô hình kinh tế lượng về độ co giãn của cầu theo giá, kết hợp cùng giao diện ứng dụng Web tương tác (Streamlit, FastAPI) và báo cáo Power BI để cung cấp các khuyến nghị điều hành trực quan cho người ra quyết định.

\subsection{Phạm vi và phương pháp tiếp cận}

Về phạm vi dữ liệu và nghiệp vụ, đề tài tập trung khai thác bộ dữ liệu giao dịch lịch sử của Global Superstore gồm 51.290 dòng chi tiết đơn hàng (tương ứng với 25.035 đơn đặt hàng độc lập của 1.590 khách hàng đối với 10.292 mã sản phẩm) trong giai đoạn 48 tháng liên tục từ tháng 01/2011 đến hết tháng 12/2014, cùng một tập dữ liệu kiểm thử nạp gia tăng gồm 500 dòng giao dịch phát sinh trong Quý 1/2015. Phạm vi phân tích bao quát toàn bộ bảy khu vực thị trường quốc tế (APAC, EU, US, LATAM, EMEA, Africa, Canada) và ba ngành hàng chính (Technology, Furniture, Office Supplies) với 17 phân nhóm sản phẩm con. Các chiều phân tích tập trung vào hiệu quả tài chính (doanh thu, lợi nhuận ròng, biên lợi nhuận), chính sách chiết khấu, chi phí và thời gian vận chuyển, hành vi mua sắm theo phân khúc khách hàng và cấu trúc giỏ hàng.

Về phương pháp tiếp cận kỹ thuật, nhóm thực hiện áp dụng quy trình phát triển hệ thống dữ liệu theo hướng mô-đun hóa và hướng cấu hình. Trong giai đoạn phân tích và thiết kế kho dữ liệu, phương pháp thiết kế bốn bước của Kimball được áp dụng chặt chẽ: xác định quy trình nghiệp vụ bán hàng làm trọng tâm, tuyên bố độ mịn (Grain) của bảng sự kiện ở mức độ chi tiết nhất là mỗi dòng sản phẩm trong một đơn hàng (\texttt{row\_id}), thiết kế các chiều phân tích chuẩn hóa và xác định các độ đo cộng được (Additive Measures). Để giải quyết bài toán làm sạch dữ liệu và đồng bộ luồng xử lý, kiến trúc Medallion ba tầng được triển khai kết hợp giữa ngôn ngữ Python (thư viện Pandas, SQLAlchemy, Psycopg2) ở khâu điều phối và sức mạnh xử lý tập hợp của ngôn ngữ SQL chuẩn PostgreSQL ở khâu biến đổi dữ liệu trong cơ sở dữ liệu (In-database ELT/ETL).

Đối với phần phát triển Hệ hỗ trợ quyết định, nhóm tiếp cận theo chu trình bốn giai đoạn của Herbert Simon gồm Thám sát (Intelligence), Thiết kế (Design), Lựa chọn (Choice) và Thực thi (Implementation). Các thuật toán dự báo san bằng mũ Holt-Winters có xu thế, thuật toán khai phá tập phổ biến Apriori và mô hình kinh tế lượng co giãn giá được cài đặt thành các gói chương trình Python độc lập trong thư mục \texttt{src/dss/}. Toàn bộ thông số cấu hình lần chạy, chỉ số đánh giá mô hình và kết quả đầu ra đều được ghi nhận bền vững xuống cơ sở dữ liệu PostgreSQL trước khi cung cấp qua giao diện lập trình ứng dụng RESTful API (FastAPI), giao diện tương tác Streamlit và bộ báo cáo quản trị Power BI Desktop. Cách tiếp cận này đảm bảo tính tách biệt rõ ràng giữa tầng lưu trữ dữ liệu, tầng xử lý thuật toán và tầng trình diễn, giúp hệ thống dễ dàng kiểm thử, bảo trì và mở rộng.
